% Sección de procedimientos de seguridad para localización y recuperación de una fuente radiactiva
\section*{Procedimientos de Localización y Recuperación}

% Punto 4: Proceso de muestreo con equipos de seguridad
\textbf{4.} Con ayuda de los equipos de seguridad se debe realizar un muestreo de la zona para tratar de localizar la fuente radiactiva. Para esto se deben formar grupos de al menos dos personas por cada equipo para asegurar que se cuente. El rastreo se debe realizar con mucho cuidado y sin escandalizar a la población ya que esto sería contraproducente.

% Punto 5: Notificación a autoridades en caso de no localización
\textbf{5.} En el supuesto caso de que no se localice la fuente radiactiva en el lapso de unas cuantas horas, se dará aviso a las autoridades locales de lo ocurrido, solicitando su ayuda para la recuperación de la fuente.

% Punto 6: Procedimientos en caso de localización fuera de área segura
\textbf{6.} En caso de que se localice la fuente y esta se encuentre fuera de su contenedor se procederá con toda la cautela y coordinación con los encargados del área segura. Si acorde a nuestros cálculos corresponde a la zona restringida y se deben realizar los cálculos de la posible exposición debida la recuperación para que no solo una persona reciba toda la dosis. Una vez hecho esto se debe simular la recuperación para mejorar los tiempos y que esta sea lo más efeciente posible. Ya teniendo la seguridad de que la recuperación se realizará de la manera más óptima se debe intentar la recuperación.

% Punto 7: Informe detallado a la Comisión Nacional de Seguridad
\textbf{7.} En su oportunidad, y en compañía del encargado de seguridad radiológica se informará detalladamente a la Comisión Nacional de Seguridad Nuclear y Salvaguardas de los hechos ocurridos para que tome conocimiento.

% Punto 8: Verificación del estado de la fuente una vez recuperada
\textbf{8.} Una vez recuperada la Fuente se le practicará una Prueba de Fuga al Contenedor y se tomarán niveles de Radiación para verificar que se encuentre en buen estado.

% Punto 9: Informe pormenorizado de los eventos y pruebas realizadas
\textbf{9.} Se redactará un informe pormenorizado de lo ocurrido y se enviará a la C.N.S.N.S., junto con una copia de los Informes de Pruebas de Fuga y de Verificación de Niveles de Radiación realizados.